\ifdefined\anonymousbuild
\else
  \def\anonymousbuild{1}
\fi

\documentclass[conference]{IEEEtran}

\usepackage[T1]{fontenc}
\usepackage[utf8]{inputenc}
\usepackage{amsmath}
\usepackage{booktabs}
\usepackage{cite}
\usepackage{hyperref}
\usepackage{xcolor}
\usepackage{listings}

\newif\ifanonymous
\ifnum\anonymousbuild=1
  \anonymoustrue
\else
  \anonymousfalse
\fi

\input{generated/metrics.tex}

\newcommand{\toolname}{Ethereum Verification Benchmark}
\newcommand{\lean}{Lean 4}
\newcommand{\repoanon}{the public benchmark artifact}
\ifanonymous
  \newcommand{\artifactrepo}{\repoanon}
\else
  \newcommand{\artifactrepo}{\url{https://github.com/lfglabs-dev/verity-benchmark}}
\fi

\lstdefinelanguage{LeanFour}{
  morekeywords={theorem,def,by,have,show,exact,simp,omega,import,namespace,end,structure,axiom,sorry,admit},
  sensitive=true,
  morecomment=[l]{--},
  morecomment=[n]{/-}{-/},
  morestring=[b]",
}
\lstset{
  language=LeanFour,
  basicstyle=\scriptsize\ttfamily,
  keywordstyle=\bfseries,
  breaklines=true,
  frame=single,
  columns=fullflexible,
  keepspaces=true
}

\title{Ethereum Verification Benchmark: An Open Suite for Evaluating AI Agents at Generating Lean 4 Proofs of Smart Contract Correctness}

\ifanonymous
\author{\IEEEauthorblockN{Anonymous Authors}}
\else
\author{\IEEEauthorblockN{Thomas Marchand}
\IEEEauthorblockA{Independent Researcher\\
Lisbon, Portugal\\
research@thomas.md}}
\fi

\begin{document}
\maketitle

\begin{abstract}
Formal verification has become a practical component of high-value smart
contract assurance, but it remains limited by the cost of writing mechanized
proofs. Recent AI agents can edit code, invoke tools, and repair compiler
errors, suggesting a path toward AI-assisted formal verification. Yet the field
lacks a benchmark that asks agents to produce complete, machine-checkable
proofs for real contract obligations rather than informal arguments or unit
tests. We present \toolname, an open evaluation suite for measuring whether AI
agents can generate \lean{} proofs of Ethereum smart contract correctness.
Built on Verity's shallow embedding of contract semantics in \lean{}, the
benchmark contains \BenchmarkTaskCount{} proof tasks across
\BenchmarkCaseCount{} deployed or ecosystem-relevant Ethereum cases. Each task
fixes the contract, specification, theorem statement, allowed files, and
verifier pipeline; hidden reference proofs are withheld from agents. The
harness rejects \texttt{sorry}, \texttt{admit}, new axioms, theorem-statement
edits, and imports unavailable in the isolated workspace. Version manifests
bind scores to task, harness, and environment fingerprints, enabling
reproducible reruns and leaderboard comparisons. We describe the benchmark
design, its v0.1 composition, and preliminary complete-suite results for three
frontier coding agents: GLM 5.2 verifies 53 of 135 tasks, MiniMax M3 verifies
38, and Kimi K2.7 verifies 31. These results indicate meaningful progress but
also show that proof generation for contract verification remains far from
solved.
\end{abstract}

\begin{IEEEkeywords}
smart contracts, formal verification, Lean 4, theorem proving, benchmarks,
AI agents, Ethereum
\end{IEEEkeywords}

\section{Introduction}

Formal verification of smart contracts has matured. Tools such as the Certora
Prover, KEVM, and K-framework based analyses are used on production protocols,
and property-based testing frameworks such as Foundry have made specification
writing part of everyday Ethereum engineering practice~\cite{certora,kevm,kframework,foundry}.
The bottleneck is increasingly not the kernel of the prover, but the human who
must translate protocol intent into precise invariants and then write the proof.
For high-value decentralized finance protocols, this labor can take weeks or
months and requires rare expertise at the intersection of protocol design,
integer arithmetic, EVM semantics, and interactive theorem proving.

AI agents change the shape of this bottleneck. Modern agents can inspect a
workspace, call a compiler, read error messages, and modify proof scripts. In
mathematics and programming, benchmarks such as miniF2F, LeanDojo, HumanEval,
and SWE-bench have made progress measurable~\cite{minif2f,leandojo,humaneval,swebench}.
However, these benchmarks do not measure the target capability needed by
Ethereum verification teams: generating complete \lean{} proofs for
state-transition correctness properties of real smart contracts under a
machine-checked verifier.

\toolname{} is designed to fill that gap. It evaluates agents on proof-only
tasks derived from Ethereum protocols and ecosystem contracts. Each task gives
the agent a fixed contract model, a fixed formal specification, and one editable
\lean{} proof file. The agent succeeds only if its submitted proof is accepted
by an independent verifier that checks the theorem statement, rejects proof
placeholders, enforces a trusted axiom ledger, and rebuilds the proof in an
isolated workspace where reference proofs are unavailable.

The benchmark is complementary to existing verification tools. Certora, KEVM,
K, Foundry, and related systems provide languages, engines, and workflows for
checking smart contract properties. \toolname{} instead evaluates the
proof-producing agent: can it complete a formal proof obligation once the
contract semantics and property are fixed? It is also distinct from general
software engineering and theorem-proving benchmarks because its tasks combine
uint256 arithmetic, storage updates, state-transition semantics, and protocol
invariants drawn from deployed Ethereum systems.

This paper makes three contributions. First, it defines a benchmark protocol for
AI-generated smart contract proofs, including task isolation, theorem-statement
checking, no-placeholder rules, and CI-enforced axiom accounting. Second, it
describes the v0.1 suite: \BenchmarkTaskCount{} tasks over
\BenchmarkCaseCount{} cases and \BenchmarkFamilyCount{} protocol families, with
coverage strongest in accounting, solvency, storage effects, and linked-list
ownership structures. Third, it reports preliminary complete-suite results for
GLM 5.2, MiniMax M3, and Kimi K2.7. The best of these solves 39\% of tasks,
suggesting that AI-assisted verification is already useful for some local proof
obligations, while harder invariants and cross-state reasoning remain open.

\section{Related Work}

\textbf{Smart contract verification.}
The K framework provides executable semantics and program analysis
infrastructure, including KEVM for Ethereum bytecode~\cite{kframework,kevm}.
Runtime Verification applies these semantics to blockchain systems and
production audits~\cite{runtimeverification}. Certora provides a widely used
specification and verification workflow for smart contracts~\cite{certora}.
Foundry brings fast Solidity testing, fuzzing, and invariant tests into the
developer workflow~\cite{foundry}. These systems check contract behavior. Our
benchmark checks whether agents can synthesize the machine-checked proof
artifacts required by a \lean{}-based verification workflow.

\textbf{AI coding benchmarks.}
HumanEval measures function-level code synthesis from natural language
prompts~\cite{humaneval}. SWE-bench evaluates agents on issue-resolution tasks
in real software repositories~\cite{swebench}. These benchmarks are important
for coding ability, but they do not require a formal proof object accepted by a
theorem prover, nor do they encode EVM-specific semantic obligations.

\textbf{LLMs and theorem proving.}
miniF2F, LeanDojo, and later theorem-proving systems including DeepSeek-Prover
evaluate or improve proof search in Lean and related environments~\cite{minif2f,leandojo,deepseekprover}.
Their tasks are primarily mathematical. \toolname{} targets a different
distribution: contract state machines, storage-frame properties, arithmetic
rounding, authorization, solvency, and protocol-specific accounting.

\section{Background: Verity and Lean Proof Tasks}

Verity is a \lean{} framework for writing smart contracts, specifications, and
proofs against a shallow embedding of contract semantics. In a shallow
embedding, contract operations are represented using host-language definitions
and propositions rather than as uninterpreted syntax alone. This gives proof
authors direct access to \lean{}'s theorem prover while preserving a semantic
model of storage, calls, arithmetic, and state transitions.

A benchmark task fixes three ingredients. The contract model defines the
operation being verified. The specification states the intended postcondition or
invariant. The target theorem connects the implementation to the specification.
The editable proof file initially contains the theorem statement and an empty
proof body; hidden reference proofs exist only for benchmark validation and are
not present in the agent workspace.

\begin{lstlisting}[caption={Schematic benchmark task shape.},label={lst:task}]
theorem deposit_sets_total_assets
  (s : VaultState) (assets : UInt256)
  (h : assets > 0) :
  spec_total_assets (deposit s assets h) =
    s.totalAssets + assets := by
  -- agent must complete this proof
\end{lstlisting}

The benchmark does not claim that every Ethereum behavior is fully modeled.
Instead, each case declares the formal surface it exercises, and the suite makes
coverage limitations explicit.

\section{Benchmark Design}

\subsection{Task Contract}

Every task is a proof-only problem. Agents may edit only designated proof files.
They may inspect public contract and specification files, generic helper
libraries, and a generated task summary. They may not see the hidden reference
proof, modify the theorem statement, add axioms, or alter the verifier. This
keeps scores focused on proof synthesis rather than benchmark adaptation.

\subsection{Axiom Ledger and Placeholder Rules}

The harness rejects \texttt{sorry}, \texttt{admit}, and new \texttt{axiom}
declarations in submitted proof files. Existing trusted assumptions are tracked
through a repository-level axiom ledger. CI enforces that the ledger remains
explicit and that proof files cannot silently expand the trusted computing base.
This is essential for a benchmark where otherwise an agent could ``solve'' a
task by assuming the target theorem.

\subsection{Theorem-Statement Integrity}

A proof attempt is not successful merely because Lean compiles. The verifier
checks that the theorem statement remains byte-for-byte consistent with the
task interface, modulo permitted proof-body edits. This prevents agents from
weakening preconditions, changing postconditions, or proving a related but
easier theorem.

\subsection{Isolated Workspaces}

Each run occurs in a fresh workspace containing public task files and generic
libraries, but not hidden proofs, private build artifacts, credentials, or
non-public imports. The final verifier then rebuilds the submitted proof in a
separate private copy. This two-step design reduces leakage from cached state
and catches imports that were not available to the agent.

\subsection{Verifier Pipeline}

The verifier checks workspace policy, theorem integrity, placeholder absence,
axiom policy, imports, and Lean compilation. A task result is counted as
verified only if all checks pass. Other outcomes are recorded as failures or
invalid artifacts, preserving information about provider failures and zero-token
runs without ranking them as complete results.

\section{Suite Composition}

Version \BenchmarkVersion{} contains \BenchmarkTaskCount{} active task
manifests over \BenchmarkCaseCount{} cases and \BenchmarkFamilyCount{} protocol
families. The tasks are drawn from Ethereum and Ethereum-adjacent protocols,
DeFi systems, token standards, bridge logic, access-control modules, and
security challenge contracts. The suite intentionally mixes simple local
storage updates with harder arithmetic and invariant-preservation tasks.

The active difficulty distribution is \BenchmarkEasyTasks{} easy,
\BenchmarkMediumTasks{} medium, \BenchmarkHardTasks{} hard, and
\BenchmarkOtherDifficultyTasks{} other. Stronger areas include accounting
conservation, local state preservation, guarded solvency, storage writes,
access-control identity, and linked-list ownership structures. Thinner areas in
v0.1 include deep reentrancy reasoning, oracle manipulation, governance and
timelock properties, temporal/liveness properties, cross-contract
compositional reasoning, cryptographic assumptions, and adversarial EVM-level
behavior.

\input{generated/tables.tex}

\section{Methodology}

\subsection{Versioned Manifests}

Scores are meaningful only relative to a fixed benchmark version. Version
\BenchmarkVersion{} records the ordered task set, Git commit, mode, budget,
and compatibility fingerprints. The task set identifier is
\BenchmarkTaskSetId{}, the harness identifier is \BenchmarkHarnessId{}, and the
environment identifier is \BenchmarkEnvironmentId{}. The source commit recorded
by the manifest is \BenchmarkGitSha{}.

Each task also has a task fingerprint over execution-relevant files and fields,
and a task-interface fingerprint over the public files and theorem interface
visible to models. This separation allows future releases to distinguish
changes that require rerunning all tasks from changes that only affect a subset.

\subsection{Rerun Planning and Leaderboards}

The rerun planner compares old and new manifests. Harness, mode, budget, or
environment changes force broad reruns. Task additions, removals, or changed
fingerprints affect only the relevant task subset when compatible. Published
leaderboards separate complete rows from partial rows, so incomplete model runs
remain visible without being ranked against complete-suite runs.

\subsection{Metrics}

The primary metric is verified tasks over total active tasks. Token counts are
reported when providers expose them, but they are secondary because provider
interfaces, forced temperatures, rate limits, and agent wrappers differ. The
v0.1 release notes specifically identify Kimi K2.7 as a provider-forced
temperature-one run, so its score should be read as a transparent preliminary
measurement rather than a perfectly controlled comparison.

\section{Preliminary Results}

Table~\ref{tab:results} reports complete-suite v0.1 runs for the three models
highlighted in this paper. GLM 5.2 verifies 53 of 135 tasks, MiniMax M3
verifies 38, and Kimi K2.7 verifies 31. These scores are high enough to show
that agents can already complete nontrivial proof obligations, but low enough
to show that smart contract proof generation is not solved by current coding
agents.

Manual inspection of successful tasks suggests that agents are strongest on
local state updates, direct arithmetic equalities, frame properties, and
specifications where the needed lemmas are close to the target theorem. Failures
often involve larger invariant chains, non-obvious arithmetic normalization,
stateful reasoning across several fields, or discovering the right helper lemma
from the library. This aligns with the suite composition: local accounting and
storage obligations are better represented and currently more tractable, while
compositional protocol reasoning remains a future target.

\section{Discussion and Limitations}

\toolname{} measures proof synthesis under a fixed formalization. It does not
measure whether an agent can discover the right specification, identify an
unknown vulnerability, or decide whether the formal model matches the deployed
bytecode. Those are separate assurance problems. The benchmark is therefore best
understood as one layer of an AI-assisted verification stack: once engineers
choose the property and formal model, can an agent discharge the proof?

The v0.1 suite is also incomplete as a representation of Ethereum security.
It has meaningful coverage of accounting, storage, access control, and solvency
properties, but less coverage of reentrancy, oracle manipulation, governance,
L2 messaging, zero-knowledge protocol integration, cryptographic primitives,
and multi-contract composition. The immediate roadmap is to expand these
families while preserving the same verifier discipline and manifest-based
comparability.

Finally, model comparisons are preliminary. The release artifacts record
per-task traces, hashes, token totals, and caveats, but providers expose
different sampling controls and execution environments. We therefore report
verified task counts conservatively and avoid claiming that v0.1 establishes a
universal model ranking.

\section{Conclusion}

AI agents are beginning to absorb part of the proof-writing workload in formal
smart contract verification, but progress needs a benchmark that is open,
reproducible, and resistant to shortcut solutions. \toolname{} provides such a
benchmark for \lean{} proofs of Ethereum contract correctness. Its v0.1 suite
contains 135 tasks from 25 real or ecosystem-relevant cases, enforces
no-placeholder proof rules, records compatibility fingerprints, and publishes
per-model artifacts. Preliminary results show measurable capability and a large
remaining gap, making the benchmark a useful target for both theorem-proving
agents and smart contract verification research.

\ifanonymous
\else
\section*{Acknowledgment}
The public benchmark repository is available at \artifactrepo. The v0.1 release
artifacts were published on June 16, 2026.
\fi

\bibliographystyle{IEEEtran}
\bibliography{references}

\end{document}
